\subsection{MULTIPLIABILITY CRITERIA}

From now on we shall assume that the normed algebra A is complete.

THEOREM 1. Let \((\pmb{x}_n)_{n\in \mathbb{N}}\) be a sequence of points in a complete normed algebra A. a) If \((\pmb{x}_n)\) is multipliable and if its product is a unit of A, then for each \(\varepsilon >0\) there exists a finite subset \(\mathbf{J}_0\) of N such that, for every finite subset L of N which does not meet \(\mathbf{J}_0\) , we have \(||e - p_{\mathrm{L}}||\leqslant \varepsilon\) .

\begin{enumerate}
\def\labelenumi{\alph{enumi})}
\setcounter{enumi}{1}
\item
  Conversely, if the sequence \((\pmb{x}_n)\) satisfies this condition, it is multipliable. Moreover, if each \(\pmb{x}_n\) is a unit, then \(\prod_{n\in \mathbb{N}}\pmb{x}_n\) is a unit.
\end{enumerate}

a) Let \(\pmb{p}\) be the product of the multipliable sequence \((\pmb{x}_n)\) , and suppose that \(\pmb{p}\) is a unit in A; then (Chapter IX, § 3, no. 7, Proposition 13) there exist \(\alpha > 0\) and \(a > 0\) such that, for all \(y \in A\) for which we have

\[
\left| \left| \boldsymbol {y} - \boldsymbol {p} \right| \right| \leqslant \alpha ,
\]

y is a unit and \(\|y^{-1}\|\leqslant a\) . By hypothesis, for every \(\varepsilon\) such that \(o<\varepsilon<\alpha\) , there is a finite subset \(H_{0}\) of N such that, for every finite subset H of N containing \(H_{0}\) , we have \(\|p_{H}-p\|\leqslant\varepsilon\) . Let \(J_{0}=[0,m]\) be an interval of N which contains \(H_{0}\) ; for each finite subset L of N which does not meet \(J_{0}\) , the integers belonging to L are all greater than those belonging to \(H_{0}\) ; hence, if \(H=H_{0}\cup L\) , we have \(p_{H}=p_{H_{0}}p_{L}\) . Now, since \(\|p_{H_{0}}-p\|\leqslant\varepsilon\leqslant\alpha\) , \(p_{H_{0}}\) is a unit, and

\[
\left| \left| e - p _ {\mathrm{H} _ {0}} ^ {- 1} p \right| \right| \leqslant \varepsilon \left| \left| p _ {\mathrm{H} _ {0}} ^ {- 1} p \right| \right| \leqslant a \varepsilon ;
\]

since \(||p_{\mathrm{H_0}}p_{\mathrm{L}} - p|| \leqslant \varepsilon\) , we deduce

\[
\left| \left| p _ {\mathrm{L}} - p _ {\mathrm{H} _ {0}} ^ {- 1} p \right| \right| \leqslant \varepsilon \left| \left| p _ {\mathrm{H} _ {0}} ^ {- 1} \right| \right| \leqslant a \varepsilon ,
\]

and finally \(||e - p_{\mathrm{L}}|| \leqslant 2a\varepsilon\) .

\begin{enumerate}
\def\labelenumi{\alph{enumi})}
\setcounter{enumi}{1}
\tightlist
\item
  Suppose that, for each \(\varepsilon > 0\) , there exists a finite subset \(J_0\) of \(\mathbf{N}\) such that, for every finite subset \(\mathbf{L}\) of \(\mathbf{N}\) which does not meet \(J_0\) , we have \(\|e - p_{\mathbf{L}}\| \leqslant \varepsilon\) . Let \(H_0 = [0, p]\) be an interval of \(\mathbf{N}\) which contains \(J_0\) ; then every finite subset \(\mathbf{H}\) of \(\mathbf{N}\) which contains \(H_0\) can be written in the form \(H_0 \cup L\) , where the integers in \(\mathbf{L}\) are all greater than those in \(H_0\) ; hence we have \(p_H = p_{H_0}p_L\) , and since \(\mathbf{L}\) does not meet \(J_0\) , \(||p_H - p_{H_0}|| \leqslant \varepsilon ||p_{H_0}||\) , and consequently \(||p_H|| = (1 + \varepsilon)||p_{H_0}||\) . If \(p_{H_0} = 0\) , the sequence \((x_n)\) is evidently multipliable and its product is 0; excluding this trivial case, there is an interval \(H_1 = [0, q]\) containing \(H_0\) and such that, for every finite subset L of N which does not meet \(H_1\) , we have \(||e - p_L|| \leqslant \varepsilon (||p_{H_0}||)^{-1}\) . As above, it follows that, for each finite subset \(H \supset H_1\) ,
\end{enumerate}

\[
\left| \left| p _ {\mathrm{H}} - p _ {\mathrm{H} _ {4}} \right| \right| \leqslant \left(\left| \left| p _ {\mathrm{H} _ {0}} \right| \right|\right) ^ {- 1} \left| \left| p _ {\mathrm{H} _ {4}} \right| \right| \varepsilon \leqslant \varepsilon (\mathrm{I} + \varepsilon).
\]

Cauchy's criterion therefore shows that \(J \to p_J\) has a limit in \(A\) with respect to the directed set \(\mathfrak{F}(N)\) .

If all the \(x_{n}\) are units, then so are all the finite partial products \(p_{J}\) ; hence for each finite subset H containing \(H_{0}\) we have

\[
\left| \left| e - p _ {\mathrm{H} _ {0}} ^ {- 1} p _ {\mathrm{H}} \right| \right| \leqslant \varepsilon ,
\]

and this shows that, in the multiplicative group G of units of A, the image under the mapping \(J \rightarrow p_{J}\) of the section filter of \(\mathfrak{F}(N)\) is a Cauchy filter base with respect to the left uniformity of G; but since G is complete (Chapter IX, § 3, no. 7, Proposition 13), the limit of the mapping \(J \rightarrow p_{J}\) belongs to G.

Remark. If \((\pmb{x}_n)\) is multipliable and its product is not a unit, the condition of Theorem 1 is not necessarily satisfied: for example, if all the \(x_n\) are equal to the same element \(x\) , where \(||x|| < 1\) , the sequence \((\pmb{x}_n)\) is multipliable and its product is 0, and for each non-empty finite subset H of N, we have \(\| p_{\mathbf{H}} \| \leqslant \| x \| < 1 \leqslant \| e \|\) .

COROLLARY I. If \((\pmb{x}_n)\) is a multipliable sequence whose product is a unit of A, then \(\lim_{n\to \infty}\pmb{x}_n = \pmb{e}\) .

COROLLARY 2. If \((\pmb{x}_n)\) is a multipliable sequence whose product is a unit of A, then every subsequence \((\pmb{x}_{n_k})_{k\in \mathbb{N}}\) of \((\pmb{x}_n)[(n_k)\) being a strictly increasing sequence of integers{]} is multipliable.

This follows immediately from the criterion of Theorem 1.

THEOREM 2. Let A be a complete normed algebra. If \((\pmb{u}_n)\) is an absolutely convergent series of elements of A, then the sequence \((\pmb{e} + \pmb{u}_n)\) is multipliable in A; and if all the elements \(\pmb{e} + \pmb{u}_n\) are units in A, then so is \(\prod_{n\in \mathbb{N}}(\pmb {e} + \pmb{u}_n)\) .

Let us apply the criterion of Theorem 1. For every finite subset L of N, we have \(p_{L} - e = \prod_{n \in L} (e + u_n) - e = \sum_{M} \left( \prod_{n \in M} u_n \right)\) , where M runs through the set of all non-empty subsets of L (linearly ordered by the induced ordering). Since \(\left\| \prod_{n\in \mathbf{M}}\pmb {u}_n\right\| \leqslant \prod_{n\in \mathbf{M}}||\pmb {u}_n||,\) we may write

\[
\left| \left| \boldsymbol {p} _ {\mathrm{L}} - \boldsymbol {e} \right| \right| \leqslant \sum_ {\mathbf {M}} \left(\prod_ {n \in \mathbf {M}} \left| \left| \boldsymbol {u} _ {n} \right|\right)\right) = \prod_ {n \in \mathbf {L}} (\mathrm{I} + \left| \left| \boldsymbol {u} _ {n} \right|\right) - \mathrm{I}.
\]

Now since the series whose general term is \(||\boldsymbol{u}_n||\) is convergent by hypothesis, the sequence \((\mathbf{I} + ||\boldsymbol{u}_n||)\) is multiplierable in \(\mathbf{R}_+^*\) (Chapter IV, § 7, no. 4, Theorem 4). Hence for each \(\varepsilon > 0\) there exists a finite subset \(J_0\) of \(\mathbf{N}\) such that, for every finite subset \(\mathbf{L}\) of \(\mathbf{N}\) which does not meet \(J_0\) , we have \(\left| \prod_{n \in \mathbf{L}} (\mathbf{I} + ||\boldsymbol{u}_n||) - \mathbf{I} \right| \leqslant \varepsilon\) ; hence the result.

COROLLARY. If the series whose general term is \(u_{n}\) is absolutely convergent, and if none of the elements \(e + u_{n}\) is a zero divisor in A, then the product \(\prod_{n \in \mathbb{N}} (e + u_{n})\) is not a zero divisor in A.

There is only a finite number of integers \(n\) such that \(||\boldsymbol{u}_n|| \geqslant 1\) . Let \(J = [0, m]\) be an interval of \(N\) containing all these integers. The product of the sequence \((\boldsymbol{e} + \boldsymbol{u}_n)\) is the product of \(p_J\) and the element \(\prod_{n > m} (\boldsymbol{e} + \boldsymbol{u}_n)\) , all of whose factors are units (Chapter IX, § 3, no. 7, Corollary to Proposition 12), and is therefore itself a unit; since \(p_J\) is the product of a finite number of non-zero divisors, it is not a zero divisor, and hence \(\prod_{n \in N} (\boldsymbol{e} + \boldsymbol{u}_n)\) is not a zero divisor.

The sufficient condition for multipliability given by Theorem 2 is not in general necessary (cf.~Exercise 6). However, it is necessary in the important case where A is an algebra of finite rank over the field R (i.e., A is finite-dimensional as a vector space over R); in particular this is the case if A is the division ring of quaternions H, or a matrix algebra \(\mathbf{M}_n(\mathbf{R})\) :

PROPOSITION 1. Let A be a normed algebra of finite rank over R. If \((\boldsymbol{e} + \boldsymbol{u}_{n})\) is a multipliable sequence in A, whose product is a unit of A, then the series whose general term is \(u_{n}\) is absolutely convergent.

From Chapter VII, § 3, no. 1, Proposition 2, there exists a number \(c > 0\) such that, for every finite family \((\pmb{x}_i)_{i \in \mathbf{I}}\) of points of \(\mathbf{A}\) , we have

\[
\sum_ {i \in \mathbf {I}} | | \boldsymbol {x} _ {i} | | \leqslant c. \sup _ {\mathrm{J} \subset \mathbf {I}} \left\| \sum_ {i \in \mathrm{J}} \boldsymbol {x} _ {i} \right\|.\tag{1}
\]

Let \((\pmb{a}_n)_{n\in \mathbb{N}}\) be an arbitrary sequence of elements of A. For each finite subset I of N, put

\[
\boldsymbol {p} _ {\mathbf {I}} = \prod_ {i \in \mathbf {I}} (\boldsymbol {e} + \boldsymbol {a} _ {i}), \quad s _ {\mathbf {I}} = \sum_ {i \in \mathbf {I}} \boldsymbol {a} _ {i}, \quad \sigma_ {\mathbf {I}} = \sum_ {i \in \mathbf {I}} | | \boldsymbol {a} _ {i} | |.
\]

LEMMA I. For each finite subset I of N, let \(\varphi(\mathrm{I}) = \sup_{\mathrm{J} \subset \mathrm{I}} ||\boldsymbol{p}_{\mathrm{J}} - \boldsymbol{e}||\) . Then for each subset J of I we have

\[
\left| \left| p _ {J} - e - s _ {J} \right| \right| \leqslant \varphi (I) \sigma_ {J}.\tag{2}
\]

The lemma is obvious if \(J\) is empty; we shall prove it by induction on the number of elements in \(J\) . Let \(J = K \cup \{j\}\) , where \(\{j\}\) is strictly larger than every \(i \in K\) . Then \(p_J = p_K(e + a_j)\) and

\[
s _ {\mathrm{J}} = s _ {\mathrm{K}} + a _ {j},
\]

so that

\[
p _ {J} - e - s _ {J} = p _ {K} - e - s _ {K} + (p _ {K} - e) a _ {j},
\]

and by the inductive hypothesis and the definition of \(\varphi(\mathbf{I})\) , we have

\[
\left| \left| p _ {J} - e - s _ {J} \right| \right| \leqslant \varphi (I) \sigma_ {K} + \varphi (I) \left| \left| a _ {j} \right| \right| = \varphi (I) \sigma_ {J},
\]

which proves the lemma.

LEMMA 2. If I is a finite subset of N such that \(\varphi(\mathbf{I}) < \mathbf{i}/c\) , then

\[
\sigma_ {\mathbf {I}} \leqslant c _ {\varphi} (\mathbf {I}) / (\mathbf {I} - c _ {\varphi} (\mathbf {I})).
\]

For since \(\sigma_{\mathbf{J}} \leqslant \sigma_{\mathbf{I}}\) for every subset \(J\) of \(I\) , we have from (2)

\[
\left| \left| s _ {J} \right| \right| \leqslant \varphi (I) \sigma_ {I} + \left| \left| p _ {J} - e \right| \right| \leqslant (I + \sigma_ {I}) \varphi (I);
\]

and since also \(\sigma_{\mathbf{I}} \leqslant c.\sup_{\mathbf{J} \in \mathbf{I}} ||s_{\mathbf{J}}||\) , from (1), it follows that

\[
\sigma_ {\mathbf {I}} \leqslant c \varphi (\mathbf {I}) (\mathbf {I} + \sigma_ {\mathbf {I}}),
\]

which leads to the result.

Now let \((e + u_{n})\) be a multipliable sequence in A, whose product is a unit; by Theorem 1 there exists a finite subset \(J_{0}\) of N such that, for each finite subset H of N which does not meet \(J_{0}\) , we have

\[
\left\| \prod_ {i \in \mathbf {H}} (e + u _ {i}) - e \right\| \leqslant 1 / 2 c.
\]

By Lemma 2, it follows that \(\sum_{i\in\mathbf{H}}||\boldsymbol{u}_{i}||-\mathrm{i}\) for every finite subset H of N which does not meet \(J_{0}\) , and hence (Chapter IV, § 7, no. 1, Theorem 1) the family \((||\boldsymbol{u}_{n}||)\) is summable in R.

